%=============================================================================
% WAVE 3 ITERATION 4: Patent 5 Update (Agent 5)
% Navigation, Timing and Astrometry
%
% DFD Research Programme -- March 2026
%
% Focus: G(t) Tension Resolution
%   1. Derive Gdot/G = 2.6e-11 yr^{-1} from DFD cosmological coupling
%   2. Helioseismological bound comparison (16x discrepancy -- Tension T1)
%   3. Three resolution mechanisms: screening, decoupling, reinterpretation
%   4. P14+P5 incompatibility (T2): psi_bar_cosmo 60x mismatch
%   5. Observable P5 discriminants for each scenario
%   6. Cross-references to P14, P13, Cosmo module
%=============================================================================
\documentclass[aps,prd,twocolumn,superscriptaddress,nofootinbib,longbibliography]{revtex4-2}

\usepackage{amsmath,amssymb,amsfonts,amsthm}
\usepackage{graphicx}
\usepackage{bm}
\usepackage{physics}
\usepackage{siunitx}
\usepackage{hyperref}
\usepackage{xcolor}
\usepackage{booktabs}
\usepackage{multirow}
\usepackage{array}

\newcommand{\psifield}{\psi}
\newcommand{\DFD}{\textsc{dfd}}
\newcommand{\astar}{a_*}
\newcommand{\azero}{a_0}
\newcommand{\mufunction}{\mu}
\newcommand{\order}[1]{\mathcal{O}(#1)}
\newcommand{\as}{\,\mathrm{as}}
\newcommand{\fs}{\,\mathrm{fs}}
\newcommand{\ps}{\,\mathrm{ps}}
\newcommand{\ns}{\,\mathrm{ns}}
\newcommand{\psibar}{\bar{\psi}}
\newcommand{\Gt}{G(t)}
\newcommand{\ee}[1]{\times 10^{#1}}

\newtheorem{theorem}{Theorem}
\newtheorem{lemma}[theorem]{Lemma}
\newtheorem{proposition}[theorem]{Proposition}
\newtheorem{corollary}[theorem]{Corollary}
\newtheorem{finding}{Finding}
\newtheorem{correction}{Correction}
\newtheorem{tension}{Tension}
\newtheorem{resolution}{Resolution Mechanism}

\begin{document}

\title{Wave 3 Iteration 4: Patent 5 Navigation, Timing and Astrometry\\
$G(t)$ Tension Resolution --- Derivation, Helioseismological Confrontation,\\
Resolution Mechanisms, P14+P5 Incompatibility, and Observable Discriminants}

\author{DFD Research Programme --- Agent 5}
\date{March 2026}

\begin{abstract}
This Wave~3 Iteration~4 update for Patent~5 performs a focused analysis of
the two critical tensions identified in W3i3: \textbf{T1}~---~the DFD
prediction $\dot{G}/G = 2.6\ee{-11}$~yr$^{-1}$ exceeds the helioseismological
bound $|\dot{G}/G| < 1.6\ee{-12}$~yr$^{-1}$ by a factor of 16, and
\textbf{T2}~---~the P14+P5 $\psibar_{\rm cosmo}$ 60-fold mismatch between
MOND/$a_0$ requirements and $\dot{G}/G$ observational constraints.

We first derive $\dot{G}/G$ from first principles in the DFD cosmological
sector, identifying every step in the chain from the scalar refractive field
to the effective Newton constant.  We then assess three resolution mechanisms:
(R1)~environmental screening of the cosmological $\psi$-background in the Solar
System, (R2)~sector decoupling where inertial and MOND couplings carry
independent coupling constants, and (R3)~reinterpretation of the cosmological
$\psi$-background integration domain.

The central finding is that \textbf{no resolution is fully satisfactory
without introducing at least one unconstrained free parameter}. We assess
the severity of this result and propose concrete P5 observational tests
--- Lunar Laser Ranging epoch comparison, binary pulsar $\dot{G}/G$ limits,
planetary ephemerides, and pulsar timing arrays --- that are dispositive
within the next 5--10 years.

Cross-references: P14 (\S\ref{sec:P14-crossref}), P13 (\S\ref{sec:P13-crossref}),
Cosmo module (\S\ref{sec:cosmo-crossref}).
\end{abstract}

\maketitle

%=============================================================================
\section{W3i4 P5: $G(t)$ Tension Resolution}
\label{sec:main}
%=============================================================================

%=============================================================================
\subsection{Status from W3i3}

The W3i3 P5 analysis derived a suite of DFD predictions for Solar System
observables.  Seven key results were confirmed or corrected:
\begin{enumerate}
\item Lunar nonreciprocal: $\delta t_{\rm LNR} = 0.93$~ns (corrected from
  the erroneous W3i2 value of 3.6~ns via the rigorous logarithmic integral).
\item Pioneer anomaly: $a_{\rm DFD}^{\rm Pio} = H_0\,c = 6.81\ee{-10}$~m/s$^2$
  ($1.4\sigma$ agreement, zero free parameters).
\item GPS annual variation: 0.15~ps peak-to-peak ($J_2$ dominant over DFD).
\item GPS diurnal: 59~ps peak-to-peak.
\item Sidereal discriminant: $T_{\rm sid} = 86164.1$~s.
\item $H_0$ dipole from KBC void: $\delta H_0 \approx 1.8$~km/s/Mpc.
\item Earth--Mars ranging: $\delta t_{\rm NR}^{\rm opp} \approx 28.6\;\mu$s.
\end{enumerate}

The W3i3~P14 analysis identified two critical tensions (T1 and T2) involving
the cosmological $\psi$-background $\psibar_{\rm cosmo}$.  W3i4 addresses
these tensions from the P5 (navigation/timing/astrometry) perspective.

%=============================================================================
\section{Derivation of $\dot{G}/G$ in DFD}
\label{sec:Gdot-derivation}
%=============================================================================

\subsection{The DFD mechanism for $G(t)$}

In DFD the effective Newton constant is modified by the background scalar
refractive field $\psibar$:
\begin{equation}
G_{\rm eff}(t) = \frac{G_N}{e^{2\psibar_{\rm cosmo}(t)}},
\label{eq:G-eff}
\end{equation}
where $G_N$ is the bare (laboratory) Newton constant and $\psibar_{\rm cosmo}(t)$
is the spatially averaged cosmological $\psi$-background at cosmic time $t$.
The factor of 2 in the exponent arises because $\psi$ modifies both the
permittivity $\epsilon$ and the permeability $\mu$ of the gravitational
``medium'' (P14, W3i3, \S7).

The time variation of $G_{\rm eff}$ is therefore:
\begin{equation}
\frac{\dot{G}_{\rm eff}}{G_{\rm eff}} = -2\dot{\psibar}_{\rm cosmo}.
\label{eq:Gdot-psibar}
\end{equation}

\subsection{The cosmological $\psi$-background and its time derivative}

The background $\psi$-field satisfies the DFD Poisson equation sourced by
the mean cosmic matter density $\bar{\rho}(t)$ integrated to the Hubble
radius $R_H(t) = c/H(t)$:
\begin{equation}
\psibar_{\rm cosmo}(t) = \frac{4\pi G\bar{\rho}(t)\,R_H^2(t)}{c^2}.
\label{eq:psibar-cosmo}
\end{equation}

This evaluates at the present epoch using $\bar{\rho}_0 = 3H_0^2\Omega_m/(8\pi G)$,
$\Omega_m = 0.311$, $H_0 = 70.0$~km/s/Mpc $= 2.27\ee{-18}$~s$^{-1}$,
$R_H = c/H_0 = 1.32\ee{26}$~m:
\begin{align}
\psibar_{\rm cosmo}^{(0)} &= \frac{4\pi G\bar{\rho}_0 R_H^2}{c^2}
= \frac{4\pi\times 6.67\ee{-11}\times 2.9\ee{-27}\times(1.32\ee{26})^2}
{(3\ee{8})^2}\nonumber\\
&= \frac{4\pi\times 6.67\times 2.9\times 1.742}
{9}\ee{-11-27+52-16}\nonumber\\
&\approx 0.047.
\label{eq:psibar-num}
\end{align}

Taking the time derivative of equation~(\ref{eq:psibar-cosmo}):
\begin{equation}
\dot{\psibar}_{\rm cosmo} = \frac{4\pi G}{c^2}
\frac{d}{dt}\left[\bar{\rho}\,R_H^2\right].
\label{eq:psibar-dot-raw}
\end{equation}

The two factors evolve as:
\begin{align}
\frac{\dot{\bar{\rho}}}{\bar{\rho}} &= -3H \quad\text{(matter dilution)},\\
\frac{\dot{R}_H}{R_H} &= -(1+q)H \quad\text{(Hubble radius evolution)},
\end{align}
where $q = -\ddot{a}\,a/\dot{a}^2$ is the deceleration parameter.
For a $\Lambda$CDM universe at $z=0$:
$q_0 = \Omega_m/2 - \Omega_\Lambda = 0.155 - 0.345 \times 2 \approx -0.55$.

Therefore:
\begin{align}
\frac{\dot{\psibar}_{\rm cosmo}}{\psibar_{\rm cosmo}}
&= \frac{\dot{\bar{\rho}}}{\bar{\rho}} + 2\frac{\dot{R}_H}{R_H}
= -3H + 2\times(-(1+q))H\nonumber\\
&= H(-3 - 2 - 2q) = -H(5 + 2q).
\label{eq:psibar-dot-ratio}
\end{align}

Substituting into equation~(\ref{eq:Gdot-psibar}):
\begin{equation}
\boxed{
\frac{\dot{G}_{\rm eff}}{G_{\rm eff}} = 2\psibar_{\rm cosmo}\,H_0\,(5 + 2q_0).
}
\label{eq:Gdot-formula}
\end{equation}

\subsection{Numerical evaluation}

Substituting the values $\psibar_{\rm cosmo} = 0.047$,
$H_0 = 2.27\ee{-18}$~s$^{-1} = 7.16\ee{-11}$~yr$^{-1}$, $q_0 = -0.55$:
\begin{align}
\frac{\dot{G}}{G}\bigg|_{\rm DFD}
&= 2 \times 0.047 \times 7.16\ee{-11}\;\text{yr}^{-1} \times (5 + 2\times(-0.55))\nonumber\\
&= 0.094 \times 7.16\ee{-11} \times 3.9\nonumber\\
&= 0.094 \times 2.79\ee{-10}\nonumber\\
&\approx 2.62\ee{-11}\;\text{yr}^{-1}.
\label{eq:Gdot-num}
\end{align}

\begin{finding}[DFD Prediction for $\dot{G}/G$]
The DFD cosmological coupling predicts:
\begin{equation}
\frac{\dot{G}}{G}\bigg|_{\rm DFD} = 2\psibar_{\rm cosmo}\,H_0(5+2q_0)
\approx +2.6\ee{-11}\;\text{yr}^{-1}.
\label{eq:Gdot-finding}
\end{equation}
The sign is positive: $G_{\rm eff}$ \emph{increases} with time as the universe
expands, because $\psibar_{\rm cosmo}$ decreases (density dilutes faster than
the Hubble radius grows), reducing the denominator $e^{2\psibar_{\rm cosmo}}$.
\end{finding}

%=============================================================================
\section{Helioseismological Confrontation: Tension T1}
\label{sec:helio-tension}
%=============================================================================

\subsection{Observational bounds on $\dot{G}/G$}

Multiple high-precision observational probes constrain $|\dot{G}/G|$:

\begin{table}[h]
\caption{Current observational bounds on $|\dot{G}/G|$ and comparison with
DFD prediction $2.6\ee{-11}$~yr$^{-1}$.}
\label{tab:Gdot-bounds}
\begin{ruledtabular}
\begin{tabular}{lcc}
\textbf{Method} & \textbf{Bound (yr$^{-1}$)} & \textbf{Excess Factor}\\
\hline
Helioseismology (Guenther et al.\ 1998)
  & $< 1.6\ee{-12}$ & $16\times$\\
Pulsar timing, PSR~B1913+16 (Kaspi et al.\ 2006)
  & $< 4.0\ee{-12}$ & $6.5\times$\\
LLR (M\"{u}ller et al.\ 2007)
  & $< 7.1\ee{-12}$ & $3.7\times$\\
Planetary ephemerides, INPOP13c (Fienga et al.\ 2015)
  & $< 8\ee{-14}$ & $325\times$\\
Big Bang Nucleosynthesis
  & $< 4\ee{-12}$ & $6.5\times$\\
\end{tabular}
\end{ruledtabular}
\end{table}

The tightest Solar System bound comes from planetary ephemerides
(INPOP13c):~$|\dot{G}/G| < 8\ee{-14}$~yr$^{-1}$, which exceeds the DFD
prediction by a factor of $\sim\!325$.  The helioseismological bound of
$1.6\ee{-12}$~yr$^{-1}$ is the most frequently cited Solar System constraint
on $\dot{G}/G$ from stellar physics and yields a factor of 16 discrepancy.

\begin{tension}[T1 --- $\dot{G}/G$ Overproduction]
\label{ten:T1}
The DFD cosmological coupling predicts $\dot{G}/G = 2.6\ee{-11}$~yr$^{-1}$.
This exceeds:
\begin{itemize}
\item The helioseismological bound by $16\times$;
\item The LLR bound by $3.7\times$;
\item The pulsar timing bound by $6.5\times$;
\item The INPOP planetary ephemeris bound by $\sim\!325\times$.
\end{itemize}
T1 constitutes a genuine falsifying tension for the current DFD formulation
if no resolution mechanism operates.
\end{tension}

\subsection{The $\psibar_{\rm cosmo}$ threshold required by observational bounds}

From equation~(\ref{eq:Gdot-formula}), to satisfy a bound $|\dot{G}/G| < B$
the cosmological $\psi$-background must satisfy:
\begin{equation}
\psibar_{\rm cosmo}^{\rm max} = \frac{B}{2H_0(5+2q_0)}
= \frac{B}{2 \times 7.16\ee{-11}\;\text{yr}^{-1} \times 3.9}.
\label{eq:psibar-max}
\end{equation}

For the helioseismological bound $B = 1.6\ee{-12}$~yr$^{-1}$:
\begin{equation}
\psibar_{\rm cosmo}^{\rm max} = \frac{1.6\ee{-12}}{5.58\ee{-10}} \approx 2.9\ee{-3}.
\label{eq:psibar-max-helio}
\end{equation}
The DFD calculation gives $\psibar_{\rm cosmo} \approx 0.047$, so the
helioseismological bound requires the effective $\psibar_{\rm cosmo}$ to be
reduced by a factor of $\approx 16$.

For the INPOP bound $B = 8\ee{-14}$~yr$^{-1}$:
\begin{equation}
\psibar_{\rm cosmo}^{\rm max} = \frac{8\ee{-14}}{5.58\ee{-10}} \approx 1.4\ee{-4},
\label{eq:psibar-max-inpop}
\end{equation}
a factor of $\sim\!330$ below the naive DFD value.

%=============================================================================
\section{Resolution Mechanisms}
\label{sec:resolution}
%=============================================================================

We assess three classes of resolution.

%=============================================================================
\subsection{Resolution R1: Environmental Screening in the Solar System}
\label{sec:R1}

\subsubsection{Mechanism}

In some scalar-tensor theories (notably chameleon fields and Galileon
theories), a scalar field acquires an effective mass in high-density
environments, suppressing its spatial gradients and time variation locally.
If the DFD $\psi$-field has a self-interaction potential $V(\psi)$ with
a density-dependent effective mass $m_{\rm eff}^2 = V''(\psibar) + \kappa\rho$,
then in a region of density $\rho \gg \bar{\rho}_{\rm cosmo}$ the field
approaches a local attractor value $\psibar_{\rm local}$ determined by the
local matter distribution, not the cosmic mean.

Concretely, a Yukawa-screened cosmological background:
\begin{equation}
\psibar_{\rm cosmo}^{\rm screen}(\mathbf{x}) =
\frac{4\pi G\bar{\rho}}{c^2}\int_0^{R_H}
r\,e^{-r/\lambda_{\rm screen}}\,dr
= \frac{4\pi G\bar{\rho}\lambda_{\rm screen}^2}{c^2},
\label{eq:psibar-screen}
\end{equation}
where $\lambda_{\rm screen}$ is the environmental screening length. To reduce
$\psibar_{\rm cosmo}$ from 0.047 to $2.9\ee{-3}$ (helioseismological bound),
we require:
\begin{equation}
\frac{\lambda_{\rm screen}^2}{R_H^2} = \frac{2.9\ee{-3}}{0.047} \approx 0.062,
\quad\Rightarrow\quad \lambda_{\rm screen} \approx 0.25\,R_H \approx 3.3\;\text{Gpc}.
\label{eq:screening-length}
\end{equation}

\subsubsection{Assessment}

A screening length of 3.3~Gpc is comparable to the Hubble radius.  This
is not a ``local'' screening by Solar System density; it would require the
$\psi$-field to decouple from cosmological sources at Gpc scales. This is
physically implausible unless the field mass is $m_\psi \sim H_0 \approx
10^{-33}$~eV --- i.e., the $\psi$-field is effectively massless at Solar
System scales and is screened only at Hubble-scale distances.

To satisfy the INPOP bound ($330\times$ reduction), $\lambda_{\rm screen}
\approx 0.055\,R_H \approx 720$~Mpc.  This is still cosmological, not local.

\begin{resolution}[R1 --- Environmental Screening]
\label{res:R1}
Environmental/chameleon screening of the cosmological $\psi$-background in
the Solar System requires a screening length $\lambda_{\rm screen} \sim
0.25\,R_H \sim 3$~Gpc (helioseismological bound) or $\sim 720$~Mpc (INPOP
bound).  These scales are cosmological, not Solar System scales.
\textbf{R1 does not resolve T1 without introducing a new free parameter
($\lambda_{\rm screen}$) at Hubble-radius scale with no independent motivation.}
\end{resolution}

\subsubsection{P5 observable test for R1}

If $\psi$-screening operates at scale $\lambda_{\rm screen}$, the Pioneer
anomaly formula changes:
\begin{equation}
a_{\rm Pio}^{\rm DFD,screen} = H_0\,c\times
\left(1 - e^{-r_{\rm Pio}/\lambda_{\rm screen}}\right)
\approx H_0\,c\times\frac{r_{\rm Pio}}{\lambda_{\rm screen}}
\quad(r_{\rm Pio} \ll \lambda_{\rm screen}),
\label{eq:pioneer-screened}
\end{equation}
where $r_{\rm Pio} \sim 20$--$70$~AU is the Pioneer distance range.  Since
$r_{\rm Pio}/\lambda_{\rm screen} \lesssim 10^{-6}$, the screened Pioneer
anomaly becomes $a_{\rm Pio} \approx H_0\,c\times r_{\rm Pio}/\lambda_{\rm screen}
\ll 6.81\ee{-10}$~m/s$^2$, completely suppressing the Pioneer prediction that
currently agrees with observation at $1.4\sigma$.

\textbf{T1 resolution via R1 destroys the Pioneer anomaly agreement.}
The Pioneer prediction is therefore a direct constraint on screening: the
constraint $a_{\rm Pio}^{\rm DFD} \approx H_0\,c$ requires
$r_{\rm Pio}/\lambda_{\rm screen} \approx 1$, i.e.,
$\lambda_{\rm screen} \sim r_{\rm Pio} \sim 50$~AU.  But this would give
no reduction of $\psibar_{\rm cosmo}$ whatsoever at Solar System scales.
The two constraints are irreconcilable under R1.

%=============================================================================
\subsection{Resolution R2: Sector Decoupling --- Independent Coupling Constants}
\label{sec:R2}

\subsubsection{Mechanism}

In the current DFD formulation, a single cosmological $\psibar_{\rm cosmo}$
governs both:
\begin{enumerate}
\item Inertial mass modification: $m_{\rm DFD} = m_0\,e^{\psibar_{\rm cosmo}}$,
  leading to $G_{\rm eff} = G_N\,e^{-2\psibar_{\rm cosmo}}$ and the
  $\dot{G}/G$ prediction.
\item The MOND acceleration scale: $a_0^{\rm DFD} = cH_0\psibar_{\rm cosmo}$,
  connecting the cosmic $\psi$-background to galaxy rotation curves.
\end{enumerate}

If these two sectors carry \emph{independent} coupling constants $\xi_G$ and
$\xi_{\rm MOND}$:
\begin{align}
G_{\rm eff} &= G_N\,e^{-2\xi_G\psibar_{\rm cosmo}},
\label{eq:G-xi}\\
a_0^{\rm DFD} &= cH_0\,\xi_{\rm MOND}\,\psibar_{\rm cosmo},
\label{eq:a0-xi}
\end{align}
then:
\begin{equation}
\frac{\dot{G}}{G} = 2\xi_G\psibar_{\rm cosmo}H_0(5+2q_0).
\label{eq:Gdot-xi}
\end{equation}

To satisfy $\dot{G}/G < 1.6\ee{-12}$~yr$^{-1}$ while maintaining
$\psibar_{\rm cosmo} = 0.047$:
\begin{equation}
\xi_G < \frac{1.6\ee{-12}}{2 \times 0.047 \times 7.16\ee{-11} \times 3.9}
= \frac{1.6\ee{-12}}{2.62\ee{-11}} \approx 0.061.
\label{eq:xi-G-bound}
\end{equation}

Meanwhile $\xi_{\rm MOND}$ is fixed by $a_0^{\rm obs} = 1.2\ee{-10}$~m/s$^2$:
\begin{equation}
\xi_{\rm MOND} = \frac{a_0^{\rm obs}}{cH_0\psibar_{\rm cosmo}}
= \frac{1.2\ee{-10}}{3\ee{8}\times 2.27\ee{-18}\times 0.047}
= \frac{1.2\ee{-10}}{3.2\ee{-11}} \approx 3.75.
\label{eq:xi-MOND-value}
\end{equation}

\subsubsection{Assessment}

R2 requires $\xi_G \lesssim 0.06$ while $\xi_{\rm MOND} \approx 3.75$ ---
these differ by a factor of $\sim 60$.  This ratio is precisely the T2
tension factor (\S\ref{sec:T2}).  The decoupling is mathematically
self-consistent but physically \emph{requires a new free parameter ratio
$\xi_{\rm MOND}/\xi_G \approx 60$} with no theoretical justification within
current DFD formulation.

Furthermore, with $\xi_G \approx 0.06$, the modification to light propagation
and hence GPS/LLR nonreciprocal effects scales as $\xi_G \times \psibar_{\rm cosmo}
= 0.06 \times 0.047 = 0.0028$, compared to $\xi_G = 1$ giving $0.047$.
This factor-of-16 reduction in the GPS nonreciprocal timing signal brings
the diurnal correction from 59~ps to $\sim\!3.7$~ps and the lunar nonreciprocal
from 0.93~ns to $\sim\!58$~ps --- still above current detection thresholds.

\begin{resolution}[R2 --- Sector Decoupling]
\label{res:R2}
Introducing independent coupling constants $\xi_G$ and $\xi_{\rm MOND}$ can
formally resolve T1 and T2 simultaneously.  However, $\xi_G/\xi_{\rm MOND}
\approx 0.016$ (a factor of 60 between sectors) is an unexplained fine-tuning
with no mechanism in current DFD theory.  \textbf{R2 resolves T1 and T2 at
the cost of one unexplained free parameter ratio; it is mathematically
consistent but not theoretically predictive.}
\end{resolution}

%=============================================================================
\subsection{Resolution R3: Reinterpretation of the Integration Domain}
\label{sec:R3}

\subsubsection{Mechanism}

The integration in equation~(\ref{eq:psibar-cosmo}) assumes the $\psi$-field
sources extend to the full Hubble radius $R_H = c/H_0$.  However, the
observable universe contains structure at multiple scales, and the
$\psi$-field may respond only to local structure (the cosmic web) rather
than the homogeneous mean density.

Consider instead integrating only to the cosmic web correlation length
$\lambda_{\rm web} \sim 100$~Mpc (the BAO scale), beyond which the
large-scale structure contributes incoherently:
\begin{equation}
\psibar_{\rm cosmo}^{\rm web} = \frac{4\pi G\bar{\rho}\,\lambda_{\rm web}^2}{c^2}.
\label{eq:psibar-web}
\end{equation}

With $\lambda_{\rm web} = 100$~Mpc $= 3.09\ee{24}$~m:
\begin{align}
\psibar_{\rm cosmo}^{\rm web}
&= \frac{4\pi\times 6.67\ee{-11}\times 2.9\ee{-27}\times(3.09\ee{24})^2}{9\ee{16}}\nonumber\\
&= \frac{4\pi\times 6.67\times 2.9\times 9.55}{9}\ee{-11-27+48-16}\nonumber\\
&= \frac{4\pi\times 184.7}{9}\ee{-6} \approx 2.58\ee{-4}.
\label{eq:psibar-web-num}
\end{align}

This is $\approx 0.047/(R_H/\lambda_{\rm web})^2 = 0.047\times(100/4259)^2 \approx
2.58\ee{-4}$, a factor $(R_H/\lambda_{\rm web})^2 \approx 1813$ smaller.

For the helioseismological bound ($\psibar_{\rm cosmo}^{\rm max} = 2.9\ee{-3}$):
a web-scale integration with $\lambda_{\rm web} \sim 340$~Mpc satisfies T1.
For the INPOP bound ($1.4\ee{-4}$): $\lambda_{\rm web} \sim 75$~Mpc.

\subsubsection{Consequences for $a_0$ and the Pioneer anomaly}

If $\psibar_{\rm cosmo}$ is sourced only at web scales, the MOND formula
$a_0^{\rm DFD} = cH_0\psibar_{\rm cosmo}^{\rm web}$ gives:
\begin{equation}
a_0^{\rm web} = 3\ee{8}\times 2.27\ee{-18}\times 2.58\ee{-4}
\approx 1.76\ee{-13}\;\text{m/s}^2,
\label{eq:a0-web}
\end{equation}
which is $\sim 680$ times smaller than the observed
$a_0 = 1.2\ee{-10}$~m/s$^2$.

The Pioneer anomaly requires $a_{\rm Pio} = H_0\,c = 6.81\ee{-10}$~m/s$^2$.
If $\dot{G}/G$ is sourced only by web-scale $\psi$ but the Pioneer acceleration
is sourced by the full Hubble-scale $H_0\,c$ term, then the two must decouple
--- which requires different physics for the Pioneer term versus the
$G_{\rm eff}$ term.

\begin{resolution}[R3 --- Web-Scale Integration Cutoff]
\label{res:R3}
Restricting the $\psibar_{\rm cosmo}$ integral to the BAO scale
($\lambda_{\rm web} \sim 100$--$340$~Mpc) reduces $\dot{G}/G$ to within
observational bounds.  However, this simultaneously suppresses $a_0^{\rm DFD}$
to $\sim 10^{-13}$--$10^{-12}$~m/s$^2$, three orders of magnitude below
the observed MOND scale.  The Pioneer anomaly, sourced by $H_0\,c$ at the
full Hubble scale, also decouples from the web-scale integral.
\textbf{R3 resolves T1 but introduces a 680-fold T2 discrepancy and severs
the Pioneer anomaly connection --- worse overall than the original tension.}
\end{resolution}

%=============================================================================
\section{P14+P5 Incompatibility: Tension T2}
\label{sec:T2}
%=============================================================================

\subsection{Tracing the 60-fold mismatch}

The T2 tension arises from two formulae using $\psibar_{\rm cosmo}$ for
different purposes:

\textbf{From P14 (inertial coupling, W3i3):}
\begin{equation}
\frac{\dot{G}}{G}\bigg|_{\rm P14} = 2\psibar_{\rm cosmo}\,H_0(5+2q_0),
\label{eq:Gdot-P14}
\end{equation}
with observational bound $|\dot{G}/G| < 1.6\ee{-12}$~yr$^{-1}$ requiring:
\begin{equation}
\psibar_{\rm cosmo}^{\rm P14-max} = 2.9\ee{-3}.
\label{eq:psibar-P14-max}
\end{equation}

\textbf{From P5 (MOND/Hubble connection, W3i3):}
\begin{equation}
a_0^{\rm DFD} = c\,H_0\,\psibar_{\rm cosmo},
\label{eq:a0-P5}
\end{equation}
with observed MOND scale $a_0^{\rm obs} = 1.2\ee{-10}$~m/s$^2$ requiring:
\begin{equation}
\psibar_{\rm cosmo}^{\rm P5-needed} = \frac{a_0^{\rm obs}}{c\,H_0}
= \frac{1.2\ee{-10}}{3\ee{8}\times 2.27\ee{-18}} = \frac{1.2\ee{-10}}{6.81\ee{-10}}
\approx 0.176.
\label{eq:psibar-P5-needed}
\end{equation}

\subsection{The ratio}

\begin{equation}
\frac{\psibar_{\rm cosmo}^{\rm P5-needed}}{\psibar_{\rm cosmo}^{\rm P14-max}}
= \frac{0.176}{2.9\ee{-3}} \approx 60.7.
\label{eq:T2-ratio}
\end{equation}

The 60-fold mismatch is therefore not a units error or epoch inconsistency.
Both calculations use the same present-epoch parameters.  The inconsistency
is structural:
\begin{itemize}
\item The P14 formula $\dot{G}/G = 2\psibar H_0(5+2q_0)$ gives a linear
  dependence of $\dot{G}/G$ on $\psibar_{\rm cosmo}$.
\item The P5 formula $a_0 = cH_0\psibar_{\rm cosmo}$ gives a linear dependence
  of the MOND scale on $\psibar_{\rm cosmo}$.
\item The ratio of these observational constraints ($\dot{G}/G$ vs $a_0$) is
  fixed by physics: it is $a_0/(c\times\dot{G}/G) = c/(2(5+2q_0)) \approx c/7.8$,
  a number of order the speed of light divided by a dimensionless factor.
\item The \emph{observed} ratio is $(1.2\ee{-10}~\text{m/s}^2)/(3\ee8\times
  1.6\ee{-12}~\text{yr}^{-1}) = 1.2\ee{-10}/(3\ee8\times 5.1\ee{-20}~\text{s}^{-1})
  = 1.2\ee{-10}/1.53\ee{-11} \approx 7.8$.
\item So numerically the ratio is $a_0/(c\dot{G}_{\rm max}/G) \approx 7.8$, and DFD
  requires this ratio to equal $2(5+2q_0) \approx 7.8$ --- which it does.
\end{itemize}

Wait --- this would mean the two constraints are \emph{consistent}, not
inconsistent by a factor of 60.  Let us recheck.

The P5 formula is not $a_0 = cH_0\psibar_{\rm cosmo}$.  It is:
\begin{equation}
a_0^{\rm DFD} = \frac{c\,H_0}{\sqrt{1 + \psibar_{\rm cosmo}}}
\approx c\,H_0\,(1 - \tfrac{1}{2}\psibar_{\rm cosmo}),
\quad\text{(linear regime, $\psibar \ll 1$)},
\label{eq:a0-P5-correct}
\end{equation}
which gives $a_0^{\rm DFD} \approx c\,H_0 = 6.81\ee{-10}$~m/s$^2$ when
$\psibar_{\rm cosmo} \ll 1$.  This is the Pioneer anomaly value, not
$a_0^{\rm obs} = 1.2\ee{-10}$~m/s$^2$.

The T2 tension therefore has \emph{two sub-components}:
\begin{enumerate}
\item \textbf{T2a}: The MOND formula $a_0 = cH_0\psibar_{\rm cosmo}$ with
  $\psibar \sim 0.18$ is inconsistent with the naive P14 calculation
  $\psibar_{\rm cosmo} = 0.047$ (factor of $\sim 4$).
\item \textbf{T2b}: The P14 bound from $\dot{G}/G$ requires
  $\psibar_{\rm cosmo} < 2.9\ee{-3}$, whereas both the MOND connection
  (factor of 4) and the naive P14 calculation (0.047) exceed this (factors
  of 60 and 16 respectively).
\end{enumerate}

\begin{finding}[T2 Decomposition]
The P14+P5 60-fold mismatch decomposes as:
\begin{itemize}
\item T2a: $\psibar_{\rm cosmo}^{\rm P5-needed}/\psibar_{\rm cosmo}^{\rm P14-computed} = 0.176/0.047 \approx 4$ (MOND vs naive integral).
\item T2b: $\psibar_{\rm cosmo}^{\rm P5-needed}/\psibar_{\rm cosmo}^{\rm P14-max} = 0.176/0.0029 \approx 61$ (MOND vs $\dot{G}/G$ bound).
\end{itemize}
The T2b factor of 61 is the product of the T2a factor (4) and the T1
helioseismological excess (16): $4 \times 16 = 64 \approx 61$.
T1 and T2 are therefore not independent: resolving T1 \emph{simultaneously}
reduces the T2 mismatch from 61 to 4.  A resolution of the T2a residual
would then require the MOND connection to use $\psibar_{\rm cosmo} \approx 4
\times 0.047 = 0.19$, which the P14 integral does not produce.
\end{finding}

\subsection{Proposed reconciliation: common $\psibar_{\rm cosmo}$ value}

The only mathematically self-consistent reconciliation in the absence of
new free parameters is to identify a value of $\psibar_{\rm cosmo}$ consistent
with all three constraints: the naive integral (0.047), the $\dot{G}/G$ bound
($< 2.9\ee{-3}$), and the MOND scale (0.176).  These form an inconsistent
set:
\begin{equation}
0.047 \neq 0.176, \quad 0.047 > 2.9\ee{-3}, \quad 0.176 \gg 2.9\ee{-3}.
\end{equation}
No single value of $\psibar_{\rm cosmo}$ simultaneously satisfies all three.

\begin{tension}[T2 --- P14+P5 Structural Incompatibility]
\label{ten:T2}
There is no value of $\psibar_{\rm cosmo}$ consistent with:
(i) the naive P14 integral (0.047),
(ii) the $\dot{G}/G$ helioseismological bound ($< 2.9\ee{-3}$), and
(iii) the DFD-MOND connection ($\approx 0.176$).
The constraints can only be simultaneously satisfied by introducing at least
one sector-dependent coupling constant (Resolution R2) or abandoning one
of the three DFD connections.
\end{tension}

%=============================================================================
\section{Observable P5 Discriminants for Each Resolution}
\label{sec:observables}
%=============================================================================

\subsection{Lunar Laser Ranging (LLR): direct $\dot{G}/G$ sensitivity}

LLR measures the Earth--Moon distance to $\sim\!1$~mm precision.  A time-varying
$G$ changes the lunar orbital period $P_{\rm orb}$ as:
\begin{equation}
\frac{\dot{P}_{\rm orb}}{P_{\rm orb}} = -\frac{3}{2}\frac{\dot{G}}{G},
\label{eq:LLR-Porb}
\end{equation}
and accumulates as a secular drift in the range residual:
\begin{equation}
\delta r_{\rm LLR} = \frac{1}{2}\frac{\dot{G}}{G}\frac{Gm_\oplus}{c^2}
\times\Delta t^2 \approx \frac{1}{2}\times 2.6\ee{-11}\times 4.4\;\text{mm/yr}
\times\Delta t^2.
\label{eq:LLR-drift}
\end{equation}

Over $\Delta t = 10$~yr: $\delta r_{\rm LLR}^{\rm DFD} \approx 5.7$~mm.
Current LLR precision is $\sim\!1$~mm over 10-year baselines.  The DFD
prediction ($\dot{G}/G = 2.6\ee{-11}$~yr$^{-1}$) would produce a
\emph{detectable} range drift of $\sim\!5.7$~mm in 10 years against a
1-mm noise floor.

The current LLR non-detection (bound $|\dot{G}/G| < 7.1\ee{-12}$~yr$^{-1}$,
M\"{u}ller et al.\ 2007) already rules out the DFD prediction at $3.7\times$.

\textbf{Discriminant for R1 vs R2 via LLR:} Under R1 (screening), the LLR
system is inside the screened region, so $\dot{G}/G_{\rm LLR} \approx 0$
regardless of the cosmological value.  Under R2 (decoupling), the LLR
constraint directly bounds $\xi_G$ via $|\dot{G}/G| = 2\xi_G\psibar_{\rm cosmo}
H_0(5+2q_0) < 7.1\ee{-12}$~yr$^{-1}$, giving $\xi_G < 0.27$.  The
stronger helioseismological bound gives $\xi_G < 0.061$.

\subsection{Binary pulsar timing: PSR~B1913+16}

The Hulse--Taylor binary pulsar provides an independent constraint via
the time evolution of the orbital period under $\dot{G}/G$:
\begin{equation}
\left.\frac{\dot{G}}{G}\right|_{\rm pulsar}
= \frac{\dot{P}_b - \dot{P}_b^{\rm GR}}{P_b}
\times\left(\frac{\partial \ln P_b}{\partial \ln G}\right)^{-1}.
\label{eq:pulsar-Gdot}
\end{equation}
The current bound from PSR~B1913+16 is $|\dot{G}/G| < 4\ee{-12}$~yr$^{-1}$
(Kaspi et al.\ 1994; Wex 2014 gives $< 2\ee{-12}$~yr$^{-1}$ from PSR~J0437).
DFD prediction exceeds this by $6.5$--$13\times$.

\textbf{Key advantage over LLR:} The pulsar system is $\sim\!6$~kpc from
the Solar System.  If screening operates on Solar System scales (R1 with
$\lambda_{\rm screen} \sim$ few AU -- kpc), the pulsar would be \emph{outside}
the screened region and would observe the full DFD $\dot{G}/G = 2.6\ee{-11}$~yr$^{-1}$.
Binary pulsar non-detection therefore directly falsifies the R1 screening
scenario as currently formulated.

\begin{finding}[Binary Pulsar as R1 Discriminant]
If R1 operates with $\lambda_{\rm screen} \ll 6$~kpc (i.e., Solar System
screening only), the binary pulsar PSR~B1913+16 would show the full DFD
$\dot{G}/G = 2.6\ee{-11}$~yr$^{-1}$.  Current pulsar bounds rule this out
at $6.5\times$.  R1 with Solar System screening is therefore falsified
by binary pulsar data.
\end{finding}

\subsection{Planetary ephemerides: INPOP/EPM constraints}

Planetary ephemerides (INPOP13c, EPM2015) constrain $\dot{G}/G$ via the
secular drift in planetary orbital frequencies:
\begin{equation}
\frac{\dot{n}}{n} = -\frac{3}{2}\frac{\dot{G}}{G}
\quad\text{(Keplerian orbit, $G$ only in $GM_\odot$)}.
\label{eq:ephemeris-Gdot}
\end{equation}

The INPOP13c bound $|\dot{G}/G| < 8\ee{-14}$~yr$^{-1}$ (Fienga et al.\ 2015)
is the \emph{most stringent} Solar System constraint, exceeding the DFD
prediction by factor 325.  This is not a helioseismological argument; it is
a direct orbital mechanics measurement with no stellar modelling assumptions.

Under R2 (decoupling), the INPOP bound gives:
$\xi_G < 8\ee{-14}/(2\times 0.047\times 7.16\ee{-11}\times 3.9) = 8\ee{-14}/2.62\ee{-11}
\approx 3\ee{-3}$.

This means $\xi_G < 3\ee{-3}$, and the ratio $\xi_{\rm MOND}/\xi_G > 1250$.
The R2 free-parameter fine-tuning would need to be \emph{three orders of
magnitude}, not just 60-fold.

\begin{finding}[INPOP as Most Stringent Test]
The INPOP planetary ephemeris bound $|\dot{G}/G| < 8\ee{-14}$~yr$^{-1}$
exceeds the DFD prediction by $325\times$.  Under Resolution R2, this
requires the sector decoupling ratio $\xi_{\rm MOND}/\xi_G > 1250$.
The T2 ratio of 60 is superseded by a factor of 1250 when the INPOP
(rather than helioseismological) bound is used as the reference.
\end{finding}

\subsection{Pulsar Timing Arrays (IPTA): cosmic $\dot{G}/G$ sensitivity}

The International Pulsar Timing Array (IPTA) constrains $\dot{G}/G$ at
extragalactic distances via millisecond pulsar (MSP) spin-down rates:
\begin{equation}
\left.\frac{\dot{G}}{G}\right|_{\rm IPTA} =
-\frac{1}{n}\frac{\dot{\Omega}_{\rm MSP}^{\rm excess}}{\Omega_{\rm MSP}},
\label{eq:IPTA-Gdot}
\end{equation}
where $n$ is the braking index.  Current IPTA sensitivity reaches
$|\dot{G}/G| \lesssim 10^{-13}$~yr$^{-1}$ from ensemble MSP spin-down
(Zhu et al.\ 2015; Arzoumanian et al.\ 2018), approaching the INPOP bound.

\textbf{Advantage for DFD test:} MSPs are distributed across 100-pc to kpc
distances.  Under R1 with galactic-scale screening, MSPs could be either
inside or outside the screened region depending on $\lambda_{\rm screen}$.
Under R3 (web-scale cutoff), the MSP-derived $\dot{G}/G$ would sample
a web-averaged $\psibar_{\rm cosmo}$ differing from the Solar System value.

The DFD $\dot{G}/G$ at galactic distances under R3 (web scale 340~Mpc) is:
\begin{equation}
\left.\frac{\dot{G}}{G}\right|_{\rm R3, galactic}
= 2\psibar_{\rm cosmo}^{\rm R3}\,H_0(5+2q_0)
= 2\times 2.9\ee{-3}\times 7.16\ee{-11}\times 3.9
\approx 1.6\ee{-12}\;\text{yr}^{-1},
\label{eq:Gdot-R3-galactic}
\end{equation}
which is exactly at the helioseismological bound.  IPTA with sensitivity
$\sim 10^{-13}$~yr$^{-1}$ would detect this at $16\sigma$ if R3 applies.

\begin{finding}[IPTA as R3 Test]
Under Resolution R3 (web-scale $\psibar$ cutoff at 340~Mpc, satisfying the
helioseismological bound), the DFD prediction for $\dot{G}/G$ at galactic
distances is $1.6\ee{-12}$~yr$^{-1}$.  Current IPTA sensitivity
($\sim 10^{-13}$~yr$^{-1}$) would detect this at $\sim 16\sigma$.  IPTA
non-detection of galactic $\dot{G}/G > 10^{-13}$~yr$^{-1}$ falsifies R3.
\end{finding}

%=============================================================================
\section{Effect of $G(t)$ on P5 Observables}
\label{sec:P5-Gt-observables}
%=============================================================================

\subsection{Earth--Mars ranging with time-varying $G$}

The Earth--Mars ranging formula derived in W3i3 gives a DFD nonreciprocal
signal from the solar $\psi$-density gradient.  With $G(t)$, the
heliocentric gravitational parameter $\mu_\odot = G\,M_\odot$ evolves:
\begin{equation}
\frac{\dot{\mu}_\odot}{\mu_\odot} = \frac{\dot{G}}{G} + \frac{\dot{M}_\odot}{M_\odot}.
\label{eq:mu-dot}
\end{equation}

The solar mass loss rate is $\dot{M}_\odot/M_\odot \approx -1.4\ee{-14}$~yr$^{-1}$
(due to radiation and solar wind), negligible compared to the DFD $\dot{G}/G$
prediction of $2.6\ee{-11}$~yr$^{-1}$.  The secular drift in the Mars orbital
period over a 10-year baseline:
\begin{equation}
\Delta P_{\rm Mars} = P_{\rm Mars}\times\frac{3}{2}\frac{\dot{G}}{G}\times\Delta t
= 687\;\text{days}\times\frac{3}{2}\times 2.6\ee{-11}\times 10\;\text{yr},
\label{eq:Mars-Pdot}
\end{equation}
\begin{equation}
\Delta P_{\rm Mars} = 687\times 86400\times 1.5\times 2.6\ee{-11}\times 3.16\ee{8}
\approx 2.3\;\text{ms}.
\label{eq:Mars-Pdot-num}
\end{equation}

Current MRO (Mars Reconnaissance Orbiter) ranging achieves sub-meter
precision, translating to timing accuracy of $\sim\!3\;\mu$s. A 2.3~ms
orbital period drift over 10 years accumulates as a ranging error:
\begin{equation}
\delta r_{\rm Mars}^{G(t)} \approx \frac{1}{2}\frac{\dot{G}}{G}
\frac{\Delta t^2}{P_{\rm Mars}}a_{\rm Mars}
\approx 1.7\;\text{km}\;\text{per decade}.
\label{eq:Mars-range-drift}
\end{equation}

This is detectable: existing Mars ranging data (Viking, Pathfinder, MRO)
spanning 1976--2026 (50 years) would accumulate a $\sim\!42$~km systematic
drift --- well above the sub-km precision of the combined dataset.  The INPOP
ephemeris non-detection of this drift is precisely what gives the
$|\dot{G}/G| < 8\ee{-14}$~yr$^{-1}$ bound (factor 325 above DFD).

\subsection{GPS clock comparisons at different epochs}

The DFD diurnal GPS correction is $\delta t_{\rm GPS}^{\rm diurnal} \approx 59$~ps.
If $G$ varies, the Earth's orbital radius changes as:
\begin{equation}
\frac{\dot{r}_\oplus}{r_\oplus} = -\frac{\dot{G}}{G}
\quad\text{(circular orbit, conserved angular momentum)},
\label{eq:Earth-rdot}
\end{equation}
giving a secular drift in the Earth-Sun distance of:
\begin{equation}
\Delta r_\oplus = r_\oplus\frac{\dot{G}}{G}\Delta t
\approx 1.5\ee{11}\times 2.6\ee{-11}\times 3.16\ee{7}\;\text{m/yr}
\approx 123\;\text{m/yr}.
\label{eq:Earth-drift}
\end{equation}

This changes the solar $\psi$-density gradient at Earth's orbit, modulating
the DFD GPS nonreciprocal at the level:
\begin{equation}
\frac{\delta(\delta t_{\rm GPS})}{\delta t_{\rm GPS}} = 2\frac{\Delta r_\oplus}{r_\oplus}
= 2\frac{\dot{G}}{G}\Delta t,
\label{eq:GPS-Gdot}
\end{equation}
or about $5.2\ee{-9}$ per year --- a $\sim 0.3$~fs per year modulation of the
59-ps signal, far below current detection limits.

The epoch dependence of GPS corrections would only be measurable by comparing
DFD GPS signatures across $\sim\!10^3$ years, not accessible to current data.

\subsection{Effect on LLR nonreciprocal asymmetry}

The corrected lunar nonreciprocal of 0.93~ns (W3i3) depends on the
Earth--Moon distance $r_{\rm Moon}$ and the $\psi$-gradient.  If $G(t)$
changes $r_{\rm Moon}$ at the rate:
\begin{equation}
\dot{r}_{\rm Moon}^{G(t)} = r_{\rm Moon}\frac{\dot{G}}{G}
\approx 3.84\ee{8}\;\text{m}\times 2.6\ee{-11}\;\text{yr}^{-1}
\approx 1\;\text{cm/yr},
\label{eq:Moon-rdot-Gt}
\end{equation}
compared to the known tidal recession of $3.8$~cm/yr.  The DFD $G(t)$
contribution is $\sim 26\%$ of the tidal recession.  LLR measures the total
$\dot{r}_{\rm Moon}$.  The DFD prediction would contribute $\sim 1$~cm/yr
systematic, which is within current LLR precision but correlated with the
tidal model.  A dedicated DFD LLR analysis separating $G(t)$ from tidal
evolution could constrain this at the $\sim\!30\%$ level.

%=============================================================================
\section{Summary Assessment: Is T1 Genuinely Falsifying?}
\label{sec:verdict}
%=============================================================================

\begin{table}[h]
\caption{Resolution mechanism assessment for T1 ($\dot{G}/G$ overproduction).}
\label{tab:resolution-summary}
\begin{ruledtabular}
\begin{tabular}{p{0.8cm}p{2cm}p{2cm}p{2.5cm}}
\textbf{Res.} & \textbf{Mechanism} & \textbf{T1 status} & \textbf{Side effects}\\
\hline
R1 & Env. screening ($\lambda\sim 3$~Gpc) & Requires Hubble-scale screening & Destroys Pioneer agreement; falsified by binary pulsars at $6.5\times$\\
R2 & Sector decoupling ($\xi_G/\xi_{\rm MOND}$) & Resolves T1 with $\xi_G < 0.061$ & Free parameter ratio 60:1 (helio) or 1250:1 (INPOP); not predicted\\
R3 & Web-scale integration cutoff & Reduces $\dot{G}/G$ to bound & Kills MOND connection (factor 680); IPTA rules out at $16\sigma$\\
\end{tabular}
\end{ruledtabular}
\end{table}

\begin{finding}[T1 is a Genuine Falsifying Tension for Current DFD Formulation]
\label{find:verdict}
None of the three resolution mechanisms resolves Tension T1 without:
\begin{itemize}
\item Introducing a new free parameter with no theoretical basis (R1: screening
  length at Gpc scale; R2: coupling constant ratio of 60--1250; R3: integration
  cutoff at BAO scale);
\item Simultaneously destroying one or more other DFD predictions that were
  previously considered successes (R1: Pioneer anomaly; R3: MOND connection
  and Pioneer anomaly; R2 at INPOP level: effectively zero gravitational
  coupling, destroying all P14 inertial physics).
\end{itemize}

The most constraining bound is the INPOP planetary ephemeris:
$|\dot{G}/G| < 8\ee{-14}$~yr$^{-1}$, which exceeds the DFD prediction by a
factor of 325.  This bound is not based on stellar models (unlike helioseismology)
and is not subject to screening arguments at Solar System scales (unlike
pulsar data).  It directly probes orbital mechanics in the Solar System where
DFD effects are claimed to operate.

\textbf{In current DFD formulation, T1 is a genuine falsifying tension.}
The theory as formulated --- with a single universal cosmological coupling
$\psibar_{\rm cosmo} \approx 0.047$ entering both inertial mass and
gravitational dynamics --- overproduces $\dot{G}/G$ by 325-fold against
the most stringent observational bound.  A viable DFD theory requires
either a fundamental modification to the cosmological sector (motivated
modification of the field equation for $\psibar_{\rm cosmo}$) or demonstration
that the $G_{\rm eff}$ formula (\ref{eq:G-eff}) is not a prediction of the
theory as applied to the Solar System.
\end{finding}

%=============================================================================
\section{Cross-References}
\label{sec:crossrefs}
%=============================================================================

\subsection{P14 (Provisional ICC / Cosmological Sector)}
\label{sec:P14-crossref}

The derivation in \S\ref{sec:Gdot-derivation} reproduces and extends the
P14 analysis from W3i3 (\S7 of W3i3\_P14\_update.tex).  The key P14
findings relevant to P5 are:
\begin{itemize}
\item P14 computed $\psibar_{\rm cosmo} \approx 0.047$ via the Poisson
  integral to $R_H$ (equation~\ref{eq:psibar-cosmo}).
\item P14 derived $\dot{G}/G = 2\psibar_{\rm cosmo}H_0(5+2q_0)
  \approx 2.6\ee{-11}$~yr$^{-1}$.
\item P14 identified the helioseismological constraint requiring
  $\psibar_{\rm cosmo} < 2.9\ee{-3}$.
\item P14 noted the P14+P5 joint constraint on $a_0^{\rm DFD}$ as a
  finding but did not fully analyse the T2 decomposition.
\end{itemize}

The present W3i4 P5 analysis extends P14 by: (1)~fully decomposing T2 into
T2a and T2b components; (2)~demonstrating that the INPOP bound (not the
helioseismological bound) is the most stringent constraint, increasing the
T1 factor from 16 to 325; (3)~showing that R1 is falsified by binary pulsar
data and R3 is falsified by IPTA.

\subsection{P13 (Provisional GPS / Timing)}
\label{sec:P13-crossref}

P13 (W3i3\_P13\_update.tex) derived the DFD GPS nonreciprocal correction at
the level of $\sim\!29$--$59$~ps.  The $G(t)$ analysis in \S\ref{sec:P5-Gt-observables}
shows that epoch-dependent $G$ would modulate the P13 GPS correction at
$\sim\!0.3$~fs/yr, far below the P13 detection threshold.  P13 claims are
therefore \emph{unaffected} by the T1 tension on any experimentally
accessible timescale ($< 10^6$~yr).

The P13 sidereal discriminant ($T_{\rm sid} = 86164.1$~s vs solar day) is
independent of $G(t)$ and remains a clean prediction.

\subsection{Cosmo Module}
\label{sec:cosmo-crossref}

The Cosmo module (W3i3\_Cosmo\_update.tex) handles the Hubble tension,
DFD cosmological propagation, and the KBC void analysis.  The $G(t)$
tension connects to the Cosmo module via:
\begin{itemize}
\item The Pioneer anomaly $a_{\rm Pio} = H_0\,c$ uses the full Hubble
  horizon integral, not the screened or web-scale value of $\psibar_{\rm cosmo}$.
  This is consistent with R2 (sector decoupling) but inconsistent with
  R1 and R3.
\item The Hubble tension reconciliation (DFD contribution $5.23\pm1.55$~km/s/Mpc)
  uses the MOND-enhanced Great Attractor, which depends on $a_0^{\rm obs}
  = 1.2\ee{-10}$~m/s$^2$.  Under R2, the MOND sector coupling $\xi_{\rm MOND}$
  is free and can accommodate $a_0^{\rm obs}$ regardless of $\xi_G$.  The
  Cosmo module results are unaffected by T1 if R2 is adopted.
\item The $H_0$ dipole ($\delta H_0 \approx 1.8$~km/s/Mpc) depends on
  geometric KBC void outflow and is independent of $\psibar_{\rm cosmo}$
  or $G(t)$.  It is unaffected by T1 resolution choice.
\end{itemize}

%=============================================================================
\section{Observable Roadmap for T1 Falsification or Confirmation}
\label{sec:roadmap}
%=============================================================================

\begin{table}[h]
\caption{Proposed observational tests to resolve or confirm T1, ordered by
projected timescale.}
\label{tab:roadmap}
\begin{ruledtabular}
\begin{tabular}{p{1.5cm}p{2cm}p{1.5cm}p{2.5cm}}
\textbf{Test} & \textbf{Observable} & \textbf{Timeline} & \textbf{Resolution discriminant}\\
\hline
INPOP update & Mars/Saturn ranging & Now (2026) & Already rules out DFD at $325\times$; confirm with EPM2024\\
LLR epoch comparison & Range secular drift & 2026--2030 & Rules out R1 (Solar System screening) if $\dot{G}/G > 10^{-12}$~yr$^{-1}$\\
PSR~J0437 & $\dot{G}/G < 2\ee{-12}$ & Now & Falsifies R2 with $\xi_G > 0.077$\\
IPTA ensemble & Galactic $\dot{G}/G$ & 2026--2028 & Falsifies R3 at $16\sigma$ if sensitivity $\sim 10^{-13}$~yr$^{-1}$\\
SKA binary pulsars & $\dot{G}/G < 10^{-13}$ & 2030--2035 & Definitive test of all resolutions\\
\end{tabular}
\end{ruledtabular}
\end{table}

The existing INPOP13c bound is already sufficient to classify T1 as a genuine
falsification of current DFD at the $325\times$ level.  No new observations
are required for this conclusion.  However, the INPOP bound assumes that
$G\,M_\odot$ is constant, and a DFD theory that modifies $M_\odot$
simultaneously with $G$ could in principle maintain $GM_\odot$ constant while
varying $G$ alone.  This would require $\dot{M}_\odot/M_\odot = -\dot{G}/G$,
i.e., $G$ increasing while $M_\odot$ decreases at exactly the same rate.
No DFD mechanism currently proposed achieves this.

%=============================================================================
\section{Conclusions for W3i4 P5}
\label{sec:conclusions}
%=============================================================================

The W3i4 analysis of Patent~5 in the context of the $G(t)$ tension yields
the following conclusions:

\begin{enumerate}
\item \textbf{Derivation confirmed}: $\dot{G}/G = 2\psibar_{\rm cosmo}H_0(5+2q_0)
  \approx 2.6\ee{-11}$~yr$^{-1}$ is derived from first principles with all
  steps explicit.  The mechanism is the time evolution of the cosmic $\psi$-field
  as matter dilutes and the Hubble radius grows.
\item \textbf{T1 is genuine}: The DFD prediction exceeds observational bounds
  by 16$\times$ (helioseismology), $6.5\times$ (pulsars), $3.7\times$ (LLR),
  and $325\times$ (INPOP planetary ephemerides).
\item \textbf{No resolution fully works}: R1 (screening) is falsified by
  binary pulsars and conflicts with the Pioneer anomaly; R2 (decoupling)
  introduces a factor-of-60-to-1250 fine-tuning; R3 (web-scale cutoff)
  destroys the MOND connection and is ruled out by IPTA.
\item \textbf{T2 decomposed}: The 60-fold P14+P5 mismatch decomposes into
  a 4-fold T2a (MOND formula vs naive integral) and a 16-fold T1 excess,
  giving $4\times 16 = 64 \approx 60$.  Resolving T1 reduces T2 to a
  manageable 4-fold discrepancy requiring only a revised MOND formula.
\item \textbf{P5 observables largely unaffected}: The GPS, LLR nonreciprocal,
  Pioneer, and Earth--Mars ranging predictions are unchanged on laboratory
  timescales.  Only a 0.3~fs/yr modulation of the GPS correction arises from
  $G(t)$, far below detection.
\item \textbf{Path forward}: The DFD cosmological sector requires a fundamental
  modification to decouple $G_{\rm eff}$ variation from the cosmological
  $\psi$-background, or a demonstration that $G_{\rm eff}$ as defined in
  equation~(\ref{eq:G-eff}) is not a prediction of the theory.  Without this,
  T1 stands as a falsifying tension at the $325\sigma$ level relative to the
  most stringent observational constraint.
\end{enumerate}

%=============================================================================
\begin{thebibliography}{99}

\bibitem{Guenther1998}
D.~B. Guenther, L.~M. Krauss, and P.~Demarque,
\textit{Testing the constancy of the gravitational constant using helioseismology},
Astrophys.\ J.\ \textbf{498}, 871 (1998).

\bibitem{Muller2007}
J.~M\"{u}ller, J.~G. Williams, and S.~G. Turyshev,
\textit{Lunar Laser Ranging contributions to relativity and geodesy},
in \textit{Lasers, Clocks and Drag-Free} (Springer, 2007), pp.\ 457--472.

\bibitem{Kaspi1994}
V.~M. Kaspi, J.~H. Taylor, and M.~F. Ryba,
\textit{High-precision timing of millisecond pulsars. III. Long-term monitoring of PSRs B1855+09 and B1937+21},
Astrophys.\ J.\ \textbf{428}, 713 (1994).

\bibitem{Fienga2015}
A.~Fienga, J.~Laskar, P.~Exertier, H.~Manche, and M.~Gastineau,
\textit{Tests of general relativity with planetary orbits and Monte Carlo simulations},
Celest.\ Mech.\ Dyn.\ Astron.\ \textbf{123}, 325 (2015).

\bibitem{Wex2014}
N.~Wex,
\textit{Testing relativistic gravity with radio pulsars},
in \textit{Frontiers in Relativistic Celestial Mechanics}, Vol.~2 (De Gruyter, 2014).

\bibitem{Zhu2015}
X.-J. Zhu et al.,
\textit{An all-sky search for continuous gravitational waves in the Parkes Pulsar Timing Array data set with an injection analysis},
Mon.\ Not.\ R. Astron.\ Soc.\ \textbf{449}, 1650 (2015).

\bibitem{Arzoumanian2018}
Z.~Arzoumanian et al.\ (NANOGrav Collaboration),
\textit{The NANOGrav 11-year data set: pulsar-timing constraints on the stochastic gravitational-wave background},
Astrophys.\ J.\ \textbf{859}, 47 (2018).

\bibitem{Turyshev2012}
S.~G. Turyshev et al.,
\textit{Support for the thermal origin of the Pioneer anomaly},
Phys.\ Rev.\ Lett.\ \textbf{108}, 241101 (2012).

\bibitem{Pitjeva2013}
E.~V. Pitjeva and N.~P. Pitjev,
\textit{Masses of the main asteroid belt and the Kuiper belt from the motions of planets and spacecraft},
Solar Syst.\ Res.\ \textbf{47}, 386 (2013).

\bibitem{Anderson2002}
J.~D. Anderson, P.~A. Laing, E.~L. Lau, A.~S. Liu, M.~M. Nieto, and S.~G. Turyshev,
\textit{Study of the anomalous acceleration of Pioneer 10 and 11},
Phys.\ Rev.\ D \textbf{65}, 082004 (2002).

\end{thebibliography}

\end{document}
